\subsection{Quantum Teleportation}\label{ex:quantum-teleportation}

\subsubsection{Problem Definition and Required Resources}\label{ex:problem-definition-and-required-resources}

Let's imagine Alice has a qubit:
\begin{equation*}
    \ket{\psi} = \alpha\ket{0} + \beta\ket{1}
\end{equation*}
And she wants Bob to have \textbf{exactly the same quantum state}. Alice does \textbf{not} know what $\alpha$ and $\beta$ are. The state could be anything: \ket{0}, \ket{1}, \ket{+}, \ket{-}, or any point on the Bloch sphere. So the protocol must work for \textbf{every possible qubit}.

\highspace
\textcolor{Green3}{\faIcon{question-circle} \textbf{Why can't Alice simply measure it?}} This is the first obstacle. If Alice measures the qubit she obtains only $0$ or $1$. Suppose the state:
\begin{equation*}
    \ket{\psi} = \frac{\sqrt{3}}{2}\ket{0} + \frac{1}{2}\ket{1}
\end{equation*}
After measuring, she might get \ket{0} or \ket{1}. The original amplitudes $\alpha$ and $\beta$ are gone forever. So \hl{measurement destroys the information we wanted to send.}

\highspace
\textcolor{Green3}{\faIcon{question-circle} \textbf{Why can't Alice simply tell Bob the amplitudes?}} Because she doesn't know them. Even worse, if she did know them ($\alpha$ and $\beta$), she would need to send an infinite amount of information to Bob. A qubit contains a \textbf{continuous} amount of information (a point on the Bloch sphere), not just one classical bit. So Alice cannot write in a message ``\emph{Hey Bob, use $\alpha = 0.834$ and $\beta = 0.551e^{i0.42}$.}'', because she has no way to discover those numbers from a single copy of the qubit.

\highspace
\textcolor{Green3}{\faIcon{question-circle} \textbf{Why can't she copy the qubit?}} We already know the answer. The No-Cloning Theorem (\autopageref{thm:no-cloning}) says: ``\emph{\textbf{there is no quantum gate that can copy an arbitrary unknown quantum state}}''. This is exactly why teleportation is so surprising.

\highspace
\begin{flushleft}
    \textcolor{Green3}{\faIcon{question-circle} \textbf{So what is teleportation actually trying to do? And why is this surprising?}}
\end{flushleft}
It is \textbf{not} send the particle. It is \textbf{transfer the quantum state}.

\highspace
Imagine Alice owns \emph{electron A} and Bob owns \emph{electron B}. At the end, Alice's electron \textbf{does not} have the state anymore; Bob's electron \textbf{does}. The particle never moved, only the \textbf{state} moved. This is why it is called \definition{State Teleportation}.

\highspace
It is surprising because Bob receives \ket{\psi} without Alice ever sending the qubit, the amplitudes, or even a copy of the qubit. Instead, they only exchange one entangled pair and two classical bits.

\newpage

\begin{flushleft}
    \textcolor{Green3}{\faIcon{question-circle} \textbf{The three resources required for teleportation}}
\end{flushleft}
We still don't know how to do teleportation, but first we need to understand the resources required. Teleportation is impossible if Alice starts with only one qubit. She needs three resources:
\begin{enumerate}
    \item \important{The unknown qubit}. This is the object we want to teleport.
    \begin{equation*}
        \ket{\psi} = \alpha\ket{0} + \beta\ket{1}
    \end{equation*}
    We can think of it as someone handling Alice a mysterious box and saying: ``\emph{Please make Bob end up with this exact quantum state}''. The important point is that \textbf{Alice does not know the values of $\alpha$ and $\beta$}. She only knows that the state has the general form:
    \begin{equation*}
        \alpha\ket{0} + \beta\ket{1}
    \end{equation*}
    The protocol must work for \textbf{every possible qubit}.


    \item \important{A shared Bell pair}. This is the most mysterious resource. Before teleportation even begins, Alice and Bob already share an entangled pair of qubits:
    \begin{equation*}
        \ket{\Phi^+} = \frac{1}{\sqrt{2}}\left(\ket{00} + \ket{11}\right)
    \end{equation*}
    This Bell pair contains \textbf{no information about the unknown qubit $\ket{\psi}$}. It is prepared \textbf{before} Alice even receives the unknown qubit.

    It is a sort of \textbf{quantum communication channel}. In classical communication, we might need an Ethernet cable, Wi-Fi, or optical fiber. Teleportation needs something different: \textbf{shared entanglement}. Without this Bell pair, there is \textbf{no quantum resource} connecting Alice and Bob.
    
    
    \item \important{Two classical bits}. After Alice measures, she gets one of only four outcomes: 00, 01, 10, or 11. She sends those two bits to Bob. \hl{Bob uses those two bits to perform a correction on his qubit.} This is the only classical communication in the protocol. It is important to note that \textbf{the classical bits do not contain any information about $\alpha$ and $\beta$}. They only tell Bob which correction to apply.
\end{enumerate}

\highspace
\begin{flushleft}
    \textcolor{Green3}{\faIcon{book} \textbf{Qubit Ordering Convention}}
\end{flushleft}
In the teleportation protocol, we will always order the qubits as follows:
\begin{enumerate}
    \item The unknown qubit $\ket{q_{\psi}}$ (Alice's qubit).
    \item Alice's half of the Bell pair $\ket{q_A}$ (Alice's qubit).
    \item Bob's half of the Bell pair $\ket{q_B}$ (Bob's qubit).
\end{enumerate}
This is just notation, it prevents mistakes.